% GENERATED FILE. Edit canonical lessons, then run npm run generate:books.
% canonical-source-hash: fnv1a64:4ee3980aee6b30b3
% canonical-lessons: TA-C353-catakam, TA-C353-patakam, TA-C353-nanmai, TA-C353-timai, TA-C353-palan

\chapter{Advantages and Drawbacks}
\label{ch:ta-advantages-and-drawbacks}
\hlchaptermodality{353}

\begin{chapteropening}
\textbf{By the end of this chapter:} \emph{I can weigh the good and the bad side of a choice.}

Complete the last lesson of chapter 353: I can weigh the good and the bad side of a choice.
\end{chapteropening}

\section[\texorpdfstring{\ta{சாதகம்} (cātakam)}{சாதகம் (cātakam)}]{\ta{சாதகம்} (cātakam) — an advantage}

\label{lesson:TA-C353-catakam}



\begin{quote}
Before the new one: say the Tamil for to lead to, then the Tamil for to affect.
\end{quote}

\subsection*{You'll want to know: \ta{சாதகம்}}
\textbf{\ta{சாதகம்}} — \emph{cātakam} — \textquotedblleft{}an advantage\textquotedblright{}.

\begin{quote}
\ta{இந்தத்} \ta{திட்டத்தில்} \ta{பல} \ta{சாதகங்கள்} \ta{இருக்கின்றன}.
\end{quote}

\begin{quote}
\emph{intat tiṭṭattil pala cātakaṅkaḷ irukkiṉṟaṉa.}
\end{quote}

\begin{quote}
This plan has many advantages.
\end{quote}

Plural: \ta{சாதகங்கள்}, \emph{cātakaṅkaḷ}.

\ta{சாதகமாக}, in one's favour: \ta{காற்று} \ta{நமக்குச்} \ta{சாதகமாக} \ta{இருக்கிறது}, the wind is in our favour.

\begin{cousinweb}
\textbf{\ta{சாதகம்}} (\emph{cātakam}) is borrowed from Sanskrit \emph{sādhaka}, \textquotedblleft{}effective, conclusive\textquotedblright{} (CDIAL 13333a), from \emph{sādh-}, \textquotedblleft{}to succeed\textquotedblright{}: what helps a thing succeed.
\end{cousinweb}

\begin{grammarlens}[title={What you've built}]
One more word for giving reasons and opinions.
\end{grammarlens}

\subsection*{Your turn}
\begin{itemize}
\raggedright
  \item \emph{Say it:} \emph{cātakam}
  \item \emph{Say it:} \emph{cātakam}, once more
  \item \emph{Say it:} \emph{intat tiṭṭattil pala cātakaṅkaḷ irukkiṉṟaṉa.} — This plan has many advantages.
  \item \emph{Recall:} say the Tamil for to lead to, then the Tamil for to affect, then say \emph{cātakam} again
\end{itemize}

\begin{tcolorbox}[breakable,colback=teal!4,colframe=teal!35!black,title={Before you move on}]
What does \ta{சாதகம்} mean? (\textquotedblleft{}An advantage\textquotedblright{}.) Say it once more.
\end{tcolorbox}
\section[\texorpdfstring{\ta{பாதகம்} (pātakam)}{பாதகம் (pātakam)}]{\ta{பாதகம்} (pātakam) — a drawback; a disadvantage}

\label{lesson:TA-C353-patakam}



\begin{quote}
Before the new one: say the Tamil for to affect, then the Tamil for an advantage.
\end{quote}

\subsection*{You'll want to know: \ta{பாதகம்}}
\textbf{\ta{பாதகம்}} — \emph{pātakam} — \textquotedblleft{}a drawback; a disadvantage\textquotedblright{}.

\begin{quote}
\ta{இந்த} \ta{முடிவில்} \ta{சில} \ta{பாதகங்களும்} \ta{இருக்கின்றன}.
\end{quote}

\begin{quote}
\emph{inta muṭivil cila pātakaṅkaḷum irukkiṉṟaṉa.}
\end{quote}

\begin{quote}
This decision has some drawbacks too.
\end{quote}

\ta{சாதக} \ta{பாதகங்கள்}, the pros and cons.

\begin{cousinweb}
\textbf{\ta{பாதகம்}} (\emph{pātakam}) is borrowed from Sanskrit \emph{bādhaka}, \textquotedblleft{}opposing, hindering, injuring\textquotedblright{} (Monier-Williams), and rhymes with its partner \ta{சாதகம்}.
\end{cousinweb}

\begin{grammarlens}[title={What you've built}]
One more word for giving reasons and opinions.
\end{grammarlens}

\subsection*{Your turn}
\begin{itemize}
\raggedright
  \item \emph{Say it:} \emph{pātakam}
  \item \emph{Say it:} \emph{pātakam}, once more
  \item \emph{Say it:} \emph{inta muṭivil cila pātakaṅkaḷum irukkiṉṟaṉa.} — This decision has some drawbacks too.
  \item \emph{Recall:} say the Tamil for to affect, then the Tamil for an advantage, then say \emph{pātakam} again
\end{itemize}

\begin{tcolorbox}[breakable,colback=teal!4,colframe=teal!35!black,title={Before you move on}]
What does \ta{பாதகம்} mean? (\textquotedblleft{}A drawback; a disadvantage\textquotedblright{}.) Say it once more.
\end{tcolorbox}
\section[\texorpdfstring{\ta{நன்மை} (naṉmai)}{நன்மை (naṉmai)}]{\ta{நன்மை} (naṉmai) — goodness; a benefit; a good thing}

\label{lesson:TA-C353-nanmai}



\begin{quote}
Before the new one: say the Tamil for an advantage, then the Tamil for a drawback.
\end{quote}

\subsection*{You'll want to know: \ta{நன்மை}}
\textbf{\ta{நன்மை}} — \emph{naṉmai} — \textquotedblleft{}goodness; a benefit; a good thing\textquotedblright{}.

\begin{quote}
\ta{உடற்பயிற்சியால்} \ta{பல} \ta{நன்மைகள்} \ta{கிடைக்கின்றன}.
\end{quote}

\begin{quote}
\emph{uṭaṟpayiṟciyāl pala naṉmaikaḷ kiṭaikkiṉṟaṉa.}
\end{quote}

\begin{quote}
Exercise brings many benefits.
\end{quote}

\ta{நன்மை} \ta{செய்}, to do good: \ta{நன்மை}, a benefit, with the abstract -\ta{மை}.

\begin{cousinweb}
\textbf{\ta{நன்மை}} (\emph{naṉmai}) is native Dravidian: DEDR 3610 lists Tamil \emph{nal}, \textquotedblleft{}good\textquotedblright{}, and \emph{naṉmai}, \textquotedblleft{}goodness, benefit\textquotedblright{}, with the abstract ending -\ta{மை}.
\end{cousinweb}

\begin{grammarlens}[title={What you've built}]
One more word for giving reasons and opinions.
\end{grammarlens}

\subsection*{Your turn}
\begin{itemize}
\raggedright
  \item \emph{Say it:} \emph{naṉmai}
  \item \emph{Say it:} \emph{naṉmai}, once more
  \item \emph{Say it:} \emph{uṭaṟpayiṟciyāl pala naṉmaikaḷ kiṭaikkiṉṟaṉa.} — Exercise brings many benefits.
  \item \emph{Recall:} say the Tamil for an advantage, then the Tamil for a drawback, then say \emph{naṉmai} again
\end{itemize}

\begin{tcolorbox}[breakable,colback=teal!4,colframe=teal!35!black,title={Before you move on}]
What does \ta{நன்மை} mean? (\textquotedblleft{}Goodness; a benefit; a good thing\textquotedblright{}.) Say it once more.
\end{tcolorbox}
\section[\texorpdfstring{\ta{தீமை} (tīmai)}{தீமை (tīmai)}]{\ta{தீமை} (tīmai) — an evil; harm}

\label{lesson:TA-C353-timai}



\begin{quote}
Before the new one: say the Tamil for a drawback, then the Tamil for goodness.
\end{quote}

\subsection*{You'll want to know: \ta{தீமை}}
\textbf{\ta{தீமை}} — \emph{tīmai} — \textquotedblleft{}an evil; harm\textquotedblright{}.

\begin{quote}
\ta{புகை} \ta{பிடிப்பதால்} \ta{உடலுக்குத்} \ta{தீமை} \ta{ஏற்படுகிறது}.
\end{quote}

\begin{quote}
\emph{pukai piṭippatāl uṭalukkut tīmai ēṟpaṭukiṟatu.}
\end{quote}

\begin{quote}
Smoking does harm to the body.
\end{quote}

X-\ta{க்குத்} \ta{தீமை}, harm to X. With \ta{நன்மை}: \ta{நன்மை} \ta{தீமை}, good and bad.

\begin{cousinweb}
\textbf{\ta{தீமை}} (\emph{tīmai}) is native Dravidian: DEDR 3267 lists Tamil \emph{tī}, \textquotedblleft{}evil\textquotedblright{}, \emph{tītu}, \textquotedblleft{}fault, harm\textquotedblright{}, and \emph{tīmai}, \textquotedblleft{}mischief, harm\textquotedblright{}, with the abstract ending -\ta{மை}.
\end{cousinweb}

\begin{grammarlens}[title={What you've built}]
One more word for giving reasons and opinions.
\end{grammarlens}

\subsection*{Your turn}
\begin{itemize}
\raggedright
  \item \emph{Say it:} \emph{tīmai}
  \item \emph{Say it:} \emph{tīmai}, once more
  \item \emph{Say it:} \emph{pukai piṭippatāl uṭalukkut tīmai ēṟpaṭukiṟatu.} — Smoking does harm to the body.
  \item \emph{Recall:} say the Tamil for a drawback, then the Tamil for goodness, then say \emph{tīmai} again
\end{itemize}

\begin{tcolorbox}[breakable,colback=teal!4,colframe=teal!35!black,title={Before you move on}]
What does \ta{தீமை} mean? (\textquotedblleft{}An evil; harm\textquotedblright{}.) Say it once more.
\end{tcolorbox}
\section[\texorpdfstring{\ta{பலன்} (palaṉ)}{பலன் (palaṉ)}]{\ta{பலன்} (palaṉ) — a benefit; a result}

\label{lesson:TA-C353-palan}



\begin{quote}
Before the new one: say the Tamil for goodness, then the Tamil for an evil.
\end{quote}

\subsection*{You'll want to know: \ta{பலன்}}
\textbf{\ta{பலன்}} — \emph{palaṉ} — \textquotedblleft{}a benefit; a result\textquotedblright{}.

\begin{quote}
\ta{கடின} \ta{உழைப்புக்குப்} \ta{பலன்} \ta{கிடைக்கும்}.
\end{quote}

\begin{quote}
\emph{kaṭiṉa uḻaippukkup palaṉ kiṭaikkum.}
\end{quote}

\begin{quote}
Hard work pays off.
\end{quote}

\ta{பலன்} \ta{கிடை}, to bear fruit, to pay off: the benefit is the subject.

\begin{cousinweb}
\textbf{\ta{பலன்}} (\emph{palaṉ}) is borrowed from Sanskrit \emph{phala}, \textquotedblleft{}fruit\textquotedblright{} (CDIAL 9051): the fruit of an effort. Tamil writes the Sanskrit ph as \ta{ப}.
\end{cousinweb}

\begin{grammarlens}[title={What you've built}]
One more word for giving reasons and opinions.
\end{grammarlens}

\subsection*{Your turn}
\begin{itemize}
\raggedright
  \item \emph{Say it:} \emph{palaṉ}
  \item \emph{Say it:} \emph{palaṉ}, once more
  \item \emph{Say it:} \emph{kaṭiṉa uḻaippukkup palaṉ kiṭaikkum.} — Hard work pays off.
  \item \emph{Recall:} say the Tamil for goodness, then the Tamil for an evil, then say \emph{palaṉ} again
\end{itemize}

\begin{tcolorbox}[breakable,colback=teal!4,colframe=teal!35!black,title={Before you move on}]
What does \ta{பலன்} mean? (\textquotedblleft{}A benefit; a result\textquotedblright{}.) Say it once more.
\end{tcolorbox}
